%% MoCafe — Single-run scans over albedo, asymmetry, and optical depth.
%% Author: Kwang-Il Seon
%% Written 2026-05-30 for kiseon_memo.cls.
\documentclass[english,a4paper]{kiseon_memo}
\synctex=1
\usepackage{amsmath}
\usepackage{amssymb}
\usepackage{microtype}

\newcommand{\rd}{\mathrm{d}}
\newcommand{\ee}{\mathrm{e}}
\newcommand{\obs}{\mathrm{obs}}
\newcommand{\edg}{\mathrm{e}}
\newcommand{\HG}{\Phi}
\newcommand{\src}{\mathrm{src}}
\newcommand{\half}{\tfrac{1}{2}}

\begin{document}

\title{Single-Run Scans over Albedo, Asymmetry, and Optical Depth in MoCafe}
\author{Kwang-Il Seon}
\date{2026-05-30}
\maketitle

\begin{abstract}
\noindent
MoCafe can record a peeling-off scattered image for an entire grid of dust
parameters from a \emph{single} Monte-Carlo run.  Photons are propagated with one
simulated asymmetry factor $g_0$ and one simulated optical depth $\tau_0$, while
closed-form per-photon weights---requiring no additional random numbers and no
re-tracing---fill in the results for a grid of albedos $a$, asymmetry factors
$g$, and optical depths $\tau$.  The scattered image becomes the five-dimensional
array $J(x,y,a,g,\tau)$.  The albedo/asymmetry reweighting follows
Seon (2010); the optical-depth (``polychromatic'') reweighting follows
Jonsson (2006).  This memo collects the underlying estimator, the reweighting
formulae, the proof that the forced first-scattering normalization cancels, and
the resulting separable form of the image accumulation.
\end{abstract}

\tableofcontents
\bigskip

%=============================================================================
\section{Baseline Monte-Carlo estimator}
%=============================================================================

MoCafe solves dust radiative transfer with a next-event (``peeling-off'')
estimator (Yusef-Zadeh, Morris \& White 1984): at every emission and scattering
site a deterministic estimate of the flux that would reach each observer is
added to the image, in addition to advancing the random walk.  The forward walk
uses two standard variance-reduction devices:

\begin{itemize}
\item \textbf{Forced first scattering} (Cashwell \& Everett 1959).  Let
  $\tau^{\edg}$ be the optical depth from the emission point to the edge of the
  grid along the initial direction.  The packet weight is multiplied by
  $(1-\ee^{-\tau^{\edg}})$---the probability of interacting before escaping---and
  the first interaction optical depth is drawn from the truncated exponential on
  $[0,\tau^{\edg}]$,
  \begin{equation}
    \tau_1 = -\ln\!\big[\,1 - R\,(1-\ee^{-\tau^{\edg}})\,\big],
    \qquad R\sim U(0,1),
    \label{eq:forced}
  \end{equation}
  whose probability density is
  \begin{equation}
    f_1(\tau_1) = \frac{\ee^{-\tau_1}}{1-\ee^{-\tau^{\edg}}},
    \qquad \tau_1\in[0,\tau^{\edg}].
    \label{eq:fdens}
  \end{equation}
\item \textbf{Continuous absorption (reduced weight).}  Beyond the first
  interaction, free paths are drawn from the unbiased exponential,
  $\tau_k=-\ln R$ with density $\ee^{-\tau_k}$, and the packet is never killed by
  absorption; instead its survival is carried as a weight.  The packet
  terminates only by escaping the grid.
\end{itemize}

Write the packet weight after the forced first scattering as
$w \equiv (1-\ee^{-\tau^{\edg}})$ (taking the emission weight to be unity).  Let
$P_k$ be the position of the $k$-th scattering, $r_k=|P_k-\mathbf{x}_\obs|$ the
distance to the observer, $\tau_\obs(P_k)$ the optical depth from $P_k$ to the
edge along the observer direction, and $\mu_{\obs,k}$ the cosine of the angle
between the incoming direction and the observer direction.  For a single dust
model with albedo $a$ and asymmetry $g$, the peeling-off contribution of the
$k$-th scattering to the image pixel onto which $P_k$ projects is
\begin{equation}
  \Delta J_k \;=\; \frac{w}{r_k^2}\;a^{k}\;
  \frac{\HG(\mu_{\obs,k}\,|\,g)}{4\pi}\;\ee^{-\tau_\obs(P_k)} ,
  \label{eq:peel}
\end{equation}
where the factor $a^{k}$ accounts for $k$ survival events (the packet has
scattered $k$ times by event $k$), and $\HG$ is the scattering phase function.
The direct (unscattered) image is the analogous estimate evaluated once at the
emission point.

MoCafe uses the Henyey--Greenstein (1941) phase function,
\begin{equation}
  \HG(\mu\,|\,g) \;=\; \frac{1-g^2}{\big(1+g^2-2g\mu\big)^{3/2}},
  \label{eq:hg}
\end{equation}
normalized so that $\int \HG(\mu|g)\,\rd\Omega/4\pi = 1$ (here $\mu=\cos\theta$ is
the cosine of the scattering angle).

%=============================================================================
\section{Albedo and asymmetry scan (Seon 2010)}
%=============================================================================

The goal is to obtain $J(x,y,a,g)$ on a grid $\{a_i\}\times\{g_l\}$ from one run.
The key observation is that the random walk's \emph{geometry}---the interaction
points and directions---can be sampled once and reweighted, because neither $a$
nor $g$ changes the extinction optical depth that drives the free paths.

\subsection{Asymmetry reweighting}

Sample the actual scattering cosines $\mu_j$ from the single simulated phase
function $\HG(\cdot\,|\,g_0)$.  A path realized with $g_0$ represents a path with
output asymmetry $g$ if each scattering is reweighted by the ratio of phase
functions (Seon 2010, eq.~2),
\begin{equation}
  W_g^{(k)} \;=\; \prod_{j=1}^{k}
  \frac{\HG(\mu_j\,|\,g)}{\HG(\mu_j\,|\,g_0)} .
  \label{eq:Wg}
\end{equation}
At the peel of event $k$ the scattering \emph{toward the observer} contributes
$\HG(\mu_{\obs,k}|g)$ directly, while the $k-1$ \emph{previous} scatterings are
represented by $W_g^{(k-1)}$.  The method is accurate while $|g-g_0|$ is not too
large; Seon (2010) finds $g_0\simeq0.4$--$0.5$ optimal for covering
$g=0.0$--$0.9$.

\subsection{Albedo reweighting}

The albedo controls only survival, not the extinction optical depth, so the
$a$-grid is filled by replacing the single factor $a^k$ in \eqref{eq:peel} with a
per-albedo accumulator
\begin{equation}
  A_a^{(k)} \;=\; a^{\,k} ,
  \label{eq:Aa}
\end{equation}
updated once per scattering (before the peel, so that the $k$-th peel sees
$a^{k}$).

\subsection{Four-dimensional image}

Combining \eqref{eq:peel}, \eqref{eq:Wg}, and \eqref{eq:Aa}, the $k$-th
scattering deposits, for every grid point $(a_i,g_l)$,
\begin{equation}
  J(x,y,a_i,g_l)\;\mathrel{+}=\;
  \frac{w}{r_k^2}\,\ee^{-\tau_\obs(P_k)}\;
  A_{a_i}^{(k)}\;\frac{\HG(\mu_{\obs,k}\,|\,g_l)}{4\pi}\;W_{g_l}^{(k-1)} .
  \label{eq:J4}
\end{equation}
The expensive geometric quantities---$\tau_\obs$ (a ray trace), $1/r_k^2$, and the
image pixel---do not depend on $(a,g)$ and are computed \emph{once} per event;
only the closed-form factors vary across the $n_a\times n_g$ grid.  Because no
random numbers are consumed by the reweighting, the slice at $(a,g_0)$ is
bit-for-bit identical to an ordinary run with that albedo and $g=g_0$.  The
direct image is independent of $a$ and $g$ and stays two-dimensional.

%=============================================================================
\section{Optical-depth (polychromatic) scan (Jonsson 2006)}
%=============================================================================

The optical-depth scan is the polychromatic algorithm of Jonsson (2006)
specialized to a \emph{grey} rescaling of one medium: all opacities are scaled by
a constant factor, so every optical depth along every ray scales by the same
amount.  Let the simulated reference be $\tau_0=$ \texttt{par\%taumax} and let the
target optical depths be $\{\tau_t\}$, with scaling factors
\begin{equation}
  s_t \;=\; \frac{\tau_t}{\tau_0} .
  \label{eq:st}
\end{equation}
Because the medium is uniformly rescaled, the optical depth to \emph{any} point
along a fixed ray satisfies $\tau^{(t)}=s_t\,\tau^{(\mathrm{ref})}$.

\subsection{Free-path reweighting}

Consider a segment in which the reference run draws free-path optical depth
$\tau_k$ and interacts at the physical point $P_k$.  Under the scaling $s_t$ the
same physical point carries optical depth $s_t\tau_k$, and the probability of
interacting there changes.  In the integration variable $\tau_k$, the reference
interaction density is $\ee^{-\tau_k}$, while the target density is
$\ee^{-s_t\tau_k}\,s_t$ (the factor $s_t$ is the Jacobian
$\rd(s_t\tau_k)/\rd\tau_k$).  The importance weight is their ratio
(Jonsson 2006, eq.~31),
\begin{equation}
  w_t^{(k)} \;=\;
  \frac{\ee^{-s_t\tau_k}\,s_t}{\ee^{-\tau_k}}
  \;=\; s_t\,\ee^{(1-s_t)\,\tau_k} ,
  \label{eq:wstep}
\end{equation}
and the cumulative per-photon optical-depth weight after $k$ segments is
\begin{equation}
  W_\tau^{(k)}(t) \;=\; \prod_{j=1}^{k} w_t^{(j)}
  \;=\; s_t^{\,k}\,\exp\!\Big[(1-s_t)\textstyle\sum_{j=1}^{k}\tau_j\Big] .
  \label{eq:Wtau}
\end{equation}
\emph{The forced first segment uses the identical factor} \eqref{eq:wstep}; the
$(1-\ee^{-\tau^{\edg}})$ normalization of forced scattering cancels in the
contribution ratio, as shown in Appendix~\ref{app:forced}.

\subsection{Peel attenuation and the escape decision}

The optical depth from a scattering site to the observer also scales, so the
extinction factor in the peel is evaluated per target,
\begin{equation}
  \ee^{-\tau_\obs(P_k)} \;\longrightarrow\; \ee^{-s_t\,\tau_\obs(P_k)} .
  \label{eq:atten}
\end{equation}
A useful property of the grey rescaling is that the \emph{interact-versus-escape}
decision is scale invariant: a non-forced segment draws $\tau_k=-\ln R$ and
escapes if $\tau_k$ exceeds the edge optical depth $D_k$; since the target edge
optical depth is $s_t D_k$ and the target interaction depth would be $s_t\tau_k$,
the ratio $\tau_k/D_k$ is independent of $s_t$.  Hence the entire forward walk is
driven by the reference run, and the scan only (i)~accumulates $W_\tau$ and
(ii)~applies the per-target attenuation \eqref{eq:atten}.  The random-number
stream is untouched, so the slice at $\tau_t=\tau_0$ (i.e.\ $s_t=1$) is
bit-for-bit identical to an ordinary single run.

%=============================================================================
\section{Combined five-dimensional estimator}
%=============================================================================

Stacking the albedo, asymmetry, and optical-depth reweightings, the $k$-th
scattering deposits, for every grid point $(a_i,g_l,\tau_t)$,
\begin{equation}
  J(x,y,a_i,g_l,\tau_t)\;\mathrel{+}=\;
  \underbrace{\frac{w}{r_k^2}}_{\text{base}}\;
  \underbrace{A_{a_i}^{(k)}}_{\text{albedo}}\;
  \underbrace{\frac{\HG(\mu_{\obs,k}|g_l)}{4\pi}\,W_{g_l}^{(k-1)}}_{\text{asymmetry}}\;
  \underbrace{\ee^{-s_t\tau_\obs(P_k)}\,W_\tau^{(k)}(t)}_{\text{optical depth}} .
  \label{eq:J5}
\end{equation}
The three parameter axes never mix inside the bracketed factors, so
\eqref{eq:J5} is the outer product of three vectors scaled by the common
geometric base $w/r_k^2$.  Defining
\begin{align}
  u_a(i) &= A_{a_i}^{(k)}, &
  v_g(l) &= \frac{\HG(\mu_{\obs,k}|g_l)}{4\pi}\,W_{g_l}^{(k-1)}, &
  q_\tau(t) &= \ee^{-s_t\tau_\obs(P_k)}\,W_\tau^{(k)}(t),
\end{align}
the per-event update is the rank-one tensor
$J_{ilt}\mathrel{+}=(w/r_k^2)\,u_a(i)\,v_g(l)\,q_\tau(t)$.  The cost is therefore
one ray trace and $n_a n_g n_t$ multiply--add operations per scattering per
observer, while the forward random walk is unchanged.

\paragraph{Direct (unscattered) light.}
Unlike $a$ and $g$, the optical depth changes the attenuation of the direct beam.
With $w_0$ the emission weight and $\tau_\obs^{\src}$ the source-to-observer
optical depth, the direct image is
\begin{equation}
  J_{\mathrm{dir}}(x,y,\tau_t)\;\mathrel{+}=\;
  \frac{w_0}{4\pi r^2}\,\ee^{-s_t\,\tau_\obs^{\src}} ,
  \label{eq:direct}
\end{equation}
which carries a $\tau$ axis but no $a$ or $g$ axis.  The unattenuated image
$J_{\mathrm{dir,0}}=w_0/(4\pi r^2)$ is independent of all three parameters and
stays two-dimensional.

%=============================================================================
\section{Implementation in MoCafe}
%=============================================================================

The scan state lives in \texttt{scan\_mod.f90}.  The per-photon accumulators
$A_a$ (\texttt{scan\_Aalb}), $W_g$ (\texttt{scan\_Wg}), and $W_\tau$
(\texttt{scan\_Wtau}) are reset at each photon birth.  Per scattering event:
$A_a$ and $W_\tau$ are advanced \emph{before} the peel (so the $k$-th peel sees
$a^k$ and the product of $k$ optical-depth factors), and $W_g$ is advanced
\emph{after} the peel using the actual scattering cosine (so the peel sees the
previous $k-1$ scatterings).  The free-path optical depth $\tau_k$ of the segment
that just ended is handed to the scan through the module variable
\texttt{scan\_tau\_step}, set in both photon-loop drivers in
\texttt{run\_simulation\_mod.f90} immediately before the call to the scatterer.

The switches and grids are
\begin{center}
\begin{tabular}{lll}
\hline
parameter & meaning & default \\
\hline
\texttt{par\%use\_ag\_list}  & enable the $(a,g)$ scan      & \texttt{.false.} \\
\texttt{par\%albedo\_list}   & albedo grid $\{a_i\}$        & $0.1,\dots,1.0$ \\
\texttt{par\%hgg\_list}      & asymmetry grid $\{g_l\}$     & $0.0,\dots,0.9$ \\
\texttt{par\%hgg}            & simulated $g_0$              & --- \\
\texttt{par\%use\_tau\_list} & enable the $\tau$ scan       & \texttt{.false.} \\
\texttt{par\%tau\_list}      & optical-depth grid $\{\tau_t\}$ & $\sqrt2$-spaced in $[0.5,2]\,\tau_0$ \\
\texttt{par\%taumax}         & simulated $\tau_0$           & --- \\
\hline
\end{tabular}
\end{center}
The two switches compose: the scattered image is $J(x,y)$, $J(x,y,a,g)$,
$J(x,y,\tau)$, or $J(x,y,a,g,\tau)$ for the four combinations.  When the $\tau$
scan is on without the $(a,g)$ scan, the $a$ and $g$ axes collapse to length one
(holding \texttt{par\%albedo} and \texttt{par\%hgg}).  The reference
$\tau_0$ slice is auto-inserted into \texttt{tau\_list} if absent.  The scan is
restricted to the no-Stokes Henyey--Greenstein path and (for $\tau$) to internal
sources.

%=============================================================================
\section{Accuracy and variance}
%=============================================================================

All reweightings are unbiased in expectation, but their variance grows as the
target parameters depart from the simulated values (Jonsson 2006, \S3.9).  For
the optical-depth scan, the per-segment weight \eqref{eq:wstep} is bounded for
$s_t>1$ (thicker targets) but the forced first segment's second moment,
$\propto\!\int_0^{\tau^{\edg}}\!\ee^{(1-2s_t)\tau_1}\rd\tau_1$, diverges with
$\tau^{\edg}$ once $s_t<\half$.  Practical guidance:
\begin{itemize}
\item keep \texttt{tau\_list} within a factor of $\sim2$--$3$ of
  \texttt{par\%taumax}, and choose \texttt{par\%taumax} near the (logarithmic)
  middle of the target range;
\item keep $g_0\simeq0.4$--$0.5$ (Seon 2010), and $|g-g_0|\lesssim0.5$--$0.6$;
\item always include the reference $\tau_0$ (done automatically) for a
  zero-variance, bit-identical anchor slice.
\end{itemize}

The implementation has been verified on a point source in a uniform sphere:
the $s_t=1$ scattered slice is bit-identical to a single run (and the direct
beam is exact at every $\tau_t$ since it involves no Monte-Carlo); the
off-reference scattered slices reproduce independent single-$\tau$ runs to
$\lesssim0.15\%$ in integrated flux, with residual pixel noise growing toward the
faint wings, as expected.

%=============================================================================
\appendix
\section{Cancellation of the forced-scattering normalization}
\label{app:forced}
%=============================================================================

We show that the forced first segment uses the same weight \eqref{eq:wstep} as an
ordinary segment, i.e.\ that the $(1-\ee^{-\tau^{\edg}})$ factor cancels in the
contribution ratio.

For the reference run the forced first scattering multiplies the packet weight by
$w=(1-\ee^{-\tau^{\edg}})$ and samples $\tau_1$ from the density $f_1$ of
\eqref{eq:fdens}.  Writing $g(\tau_1)$ for the (parameter-dependent) peel factors
that multiply the weight at the interaction point, the reference image
expectation per emitted photon is
\begin{equation}
  E_{\mathrm{ref}}
  = \int_0^{\tau^{\edg}}\! f_1(\tau_1)\,w\,g(\tau_1)\,\rd\tau_1
  = \int_0^{\tau^{\edg}}\! \ee^{-\tau_1}\,g(\tau_1)\,\rd\tau_1 ,
\end{equation}
where the normalization $(1-\ee^{-\tau^{\edg}})$ in $w$ cancels the same factor in
the denominator of $f_1$.

For an independent run at scaling $s$ (target optical depth $s\tau_0$), the edge
depth is $s\tau^{\edg}$ and the forced first scattering samples $\tau_1'$ from
$\ee^{-\tau_1'}/(1-\ee^{-s\tau^{\edg}})$ with weight $(1-\ee^{-s\tau^{\edg}})$.
Substituting $\tau_1'=s\tau_1$ (so that the physical interaction point, hence
$g$, coincides with the reference's) gives
\begin{equation}
  E_{\mathrm{indep}}(s)
  = \int_0^{\tau^{\edg}}\! \ee^{-s\tau_1}\,s\,g_s(\tau_1)\,\rd\tau_1 ,
\end{equation}
where $g_s$ differs from $g$ only through the per-target peel attenuation
$\ee^{-s\tau_\obs}$ already accounted for in \eqref{eq:atten}.  Reweighting the
reference samples by \eqref{eq:wstep} reproduces this exactly:
\begin{equation}
  \int_0^{\tau^{\edg}}\! \ee^{-\tau_1}\,g_s(\tau_1)\,
  \big[\,s\,\ee^{(1-s)\tau_1}\big]\,\rd\tau_1
  = \int_0^{\tau^{\edg}}\! s\,\ee^{-s\tau_1}\,g_s(\tau_1)\,\rd\tau_1
  = E_{\mathrm{indep}}(s) .
\end{equation}
The factor $(1-\ee^{-\tau^{\edg}})$ never appears, so the forced first segment is
weighted by $s\,\ee^{(1-s)\tau_1}$, identical in form to an ordinary
(non-forced) segment.  The same change of variables, applied segment by segment,
establishes the unbiasedness of \eqref{eq:Wtau} for all scattering orders.
\hfill$\blacksquare$

%=============================================================================
\begin{thebibliography}{9}
%=============================================================================
\bibitem{cashwell1959} Cashwell, E.~D., \& Everett, C.~J. 1959,
  \textit{A Practical Manual on the Monte Carlo Method for Random Walk
  Problems} (Pergamon Press).
\bibitem{hg1941} Henyey, L.~G., \& Greenstein, J.~L. 1941, ApJ, 93, 70.
\bibitem{jonsson2006} Jonsson, P. 2006, MNRAS, 372, 2.
\bibitem{seon2010} Seon, K.-I. 2010, PKAS, 25, 177.
\bibitem{yzmw1984} Yusef-Zadeh, F., Morris, M., \& White, R.~L. 1984, ApJ,
  278, 186.
\end{thebibliography}

\end{document}
